\section{Introduction}
Modern cloud infrastructures are increasingly multi-tenant, hosting sensitive applications and proprietary data alongside potentially malicious workloads. Historically, organizations relied on perimeter-based firewalls to isolate internal networks. However, this model fails when attackers breach the perimeter or when malicious insiders exploit internal trust. To address these vulnerabilities, the zero-trust network architecture (ZTNA) model has emerged, operating on the principle of ``never trust, always verify.''

In a zero-trust environment, every endpoint must verify its identity and security posture before accessing sensitive resources. A key mechanism for verifying endpoint integrity is \emph{remote attestation}, which allows a verifier to cryptographically confirm the software stack running on a remote system.

In this work, we analyze the security and performance of \emph{Nexa}, an enterprise-ready attestation platform developed by Nexa Corp. Nexa integrates hardware-enforced enclaves (such as Vertex Secure Enclaves) with the Synapse key management protocol to provide continuous authentication. While Nexa promises strong security guarantees, its performance under high network load and its resilience to cryptographic attacks have not been independently evaluated.

This paper provides the first comprehensive security analysis of Nexa. Our contributions are threefold:
\begin{enumerate}
    \item We detail the Nexa system architecture and the Synapse cryptographic exchange, highlighting potential vectors for exploit.
    \item We perform a physical performance evaluation of Nexa, measuring latency, CPU utilization, and memory footprint under varying workloads.
    \item We demonstrate that while Nexa provides robust resistance to protocol-level attacks, older enclave implementations present minor side-channel risks that require software-level mitigation.
\end{enumerate}

The remainder of this paper is structured as follows. Section~\ref{sec:related} reviews related work. Section~\ref{sec:method} outlines the Nexa method and architecture. Section~\ref{sec:experiments} details our experimental setup. Section~\ref{sec:results} presents the results of our evaluation. Section~\ref{sec:discussion} discusses security and scalability implications, and Section~\ref{sec:conclusion} concludes.
